\chapter{The Gaussian family}
\label{ch:gauss}

\section{Affine images}

\begin{resultbox}[Homework 1, the basic transform]
If $x\sim\N(\mu,\Sigma)$ and $y=Ax+b$, then
\[
y\sim\N(A\mu+b,\, A\Sigma A^\top),
\]
on the affine image of the column space. If $A\Sigma A^\top$ is singular, $y$ has no density on the ambient Euclidean space. It has a density on the affine hull of its support.
\end{resultbox}

The moment proof is two lines. $\E[y]=A\mu+b$ by linearity. The covariance is
\[
\E[(y-\E y)(y-\E y)^\top]=A\,\E[(x-\mu)(x-\mu)^\top]\,A^\top=A\Sigma A^\top.
\]
A Gaussian is determined by its mean and covariance, so the law is that Gaussian. The same calculation with $A$ a row vector is the reparameterization trick in one coordinate: if $\varepsilon\sim\N(0,1)$, then $\mu+\sigma\varepsilon\sim\N(\mu,\sigma^2)$.

A law defined by its moment generating function covers singular covariances as well: for every $s$,
\[
\E[e^{s^\top y}]=e^{s^\top b}\,\E[e^{(A^\top s)^\top x}]
=\exp\bigl(s^\top(A\mu+b)+\tfrac12 s^\top(A\Sigma A^\top)s\bigr).
\]
Jointly Gaussian random vectors are independent if and only if they are uncorrelated.

\begin{practbox}
Let $x\sim\N(1,4)$ and $y=3x-2$. Find $\E[y]$ and $\mathrm{Var}(y)$.
\end{practbox}
\begin{solbox}
Mean $3\cdot 1-2=1$. Variance $3^2\cdot 4=36$. So $y\sim\N(1,36)$.
\end{solbox}

\section{Product of Gaussian densities}

For $\Sigma\succ 0$,
\[
\N(x\given\mu,\Sigma)=(2\pi)^{-d/2}(\det\Sigma)^{-1/2}\exp\bigl(-\tfrac12(x-\mu)^\top\Sigma^{-1}(x-\mu)\bigr).
\]
The product of two such densities, as a function of $x$, has precision equal to the sum of the precisions.
\begin{resultbox}[Gaussian product identity]
\[
\N(x\given a,A)\,\N(x\given b,B)=\N(a\given b,A+B)\,\N(x\given c,C),
\]
where
\[
C=(A^{-1}+B^{-1})^{-1},\qquad
c=C(A^{-1}a+B^{-1}b),
\]
or, equivalently, the parallel-sum forms
\[
C=A(A+B)^{-1}B,\qquad
c=B(A+B)^{-1}a+A(A+B)^{-1}b.
\]
\end{resultbox}
Integrating the product in $x$ produces the constant $Z=\N(a\given b,A+B)$, because the remaining Gaussian integrates to one. If $U\sim\N(a,A)$ and $V\sim\N(b,B)$ are independent, that constant is the density of $U-V$ at $0$.

Homework~1 asks this identity because it is Bayes' rule for a Gaussian likelihood and a Gaussian prior: posterior precision equals prior precision plus likelihood precision.

\section{Conditionals without block inversion}

For a joint Gaussian, the conditional of one block given another is Gaussian. Its mean is linear in the observation, and its covariance does not depend on the observation. That is why a linear-Gaussian encoder can be written in closed form, and why a neural encoder is the nonlinear replacement of that conditional mean and variance.

One can avoid inverting the full joint covariance. Set $K=\Sigma_{yx}\Sigma_{xx}^{-1}$ (assuming $\Sigma_{xx}\succ 0$) and form the residual $r=y-\mu_y-K(x-\mu_x)$. Then $\mathrm{Cov}(r,x)=0$. Joint Gaussianity upgrades uncorrelated to independent, so
\begin{resultbox}[Gaussian conditional]
\[
y\given x=x_0\sim\N\bigl(\mu_y+K(x_0-\mu_x),\;\Sigma_{yy}-\Sigma_{yx}\Sigma_{xx}^{-1}\Sigma_{xy}\bigr).
\]
\end{resultbox}
The subtracted term is positive semidefinite, so $\mathrm{Cov}(y\given x)\preceq\Sigma_{yy}$, with equality if and only if $\Sigma_{yx}=0$, which for jointly Gaussian vectors is independence.

\section{Linear-Gaussian models and Woodbury}

Write $y=Ax+b+\varepsilon$ with $\varepsilon\sim\N(0,S)$ independent of $x\sim\N(\mu,\Sigma)$. The joint $(x,y)$ is Gaussian with
\[
\begin{pmatrix}x\\ y\end{pmatrix}
\sim\N\Bigg(
\begin{pmatrix}\mu\\ A\mu+b\end{pmatrix},
\begin{pmatrix}\Sigma & \Sigma A^\top\\ A\Sigma & S+A\Sigma A^\top\end{pmatrix}
\Bigg).
\]
Conditioning $x$ on $y$ by the residual argument above inverts only $R:=S+A\Sigma A^\top$:
\[
\mathrm{Cov}(x\given y)=\Sigma-\Sigma A^\top R^{-1}A\Sigma,\qquad
\E[x\given y]=\mu+\Sigma A^\top R^{-1}\bigl(y-(A\mu+b)\bigr).
\]
The information form adds prior and likelihood precisions:
\[
\Sigma_{x\given y}=(\Sigma^{-1}+A^\top S^{-1}A)^{-1},\qquad
\E[x\given y]=\Sigma_{x\given y}\bigl(\Sigma^{-1}\mu+A^\top S^{-1}(y-b)\bigr).
\]
\begin{resultbox}[Woodbury for the posterior covariance]
\[
(\Sigma^{-1}+A^\top S^{-1}A)^{-1}=\Sigma-\Sigma A^\top R^{-1}A\Sigma.
\]
\end{resultbox}
The matrix in front of the innovation $y-(A\mu+b)$ is $\mathrm{Cov}(x,y)\mathrm{Cov}(y)^{-1}$, the linear regression coefficient of $x$ on $y$.

\section{Mixtures}

A mixture $p(x)=\sum_{k=1}^K \pi_k\,\N(x\given\mu_k,\Sigma_k)$ with $\pi_k>0$ and $\sum_k\pi_k=1$ is nonnegative and integrates to one. Its mean is $\sum_k\pi_k\mu_k$. Its covariance is not the mixture of the covariances alone. The law of total variance adds the covariance of the component means:
\[
\mathrm{Cov}(x)=\E[\mathrm{Cov}(x\given k)]+\mathrm{Cov}(\E[x\given k]).
\]
In symbols, with $\mu=\sum_k\pi_k\mu_k$,
\begin{resultbox}[Mixture moments]
\[
\mathrm{Cov}(x)
=\underbrace{\sum_k\pi_k\Sigma_k}_{\text{within}}
+\underbrace{\sum_k\pi_k(\mu_k-\mu)(\mu_k-\mu)^\top}_{\text{between}}.
\]
\end{resultbox}
A single Gaussian cannot represent well-separated modes. A mixture can. Lecture~4, slides~3--11, uses the iris petals as the picture. The VAE notebook, part~II, trains that comparison in two dimensions.

On the line, the equal mixture of $\N(-a,1)$ and $\N(a,1)$ has mean $0$, variance $1+a^2$, and excess kurtosis $-2a^4/(1+a^2)^2<0$ for $a>0$. A nondegenerate Gaussian has excess kurtosis $0$, so the mixture is not Gaussian. Its modes bifurcate at $a=1$: one mode for $a\le 1$, two modes for $a>1$.

\section{Maximum entropy and the Pythagorean identity}

Among densities on $\R^d$ with a fixed mean and a fixed covariance, the Gaussian has the largest entropy
\[
h(p)=-\int p\log p.
\]
For $\N(\mu,\Sigma)$ the entropy is $\tfrac{d}{2}\log(2\pi e)+\tfrac12\log\det\Sigma$. Subtracting two such entropies, and accounting for the cross term that appears when the means differ, is one derivation of the multivariate KL~\eqref{eq:mvnkl} in Chapter~\ref{ch:info}.

Let $p$ be any density with mean $\mu_p$ and covariance $\Sigma_p\succ 0$, let $g=\N(\mu_p,\Sigma_p)$, and let $q$ be any Gaussian. The function $\log(g/q)$ is a polynomial of degree at most $2$. Any such polynomial has the same expectation under $p$ and under $g$, because $p$ and $g$ share moments of order $0,1,2$. Therefore
\begin{resultbox}[Pythagorean identity for Gaussian projection]
\[
\KL(p\Vert q)=\KL(p\Vert g)+\KL(g\Vert q).
\]
\end{resultbox}
Since $\KL(g\Vert q)\ge 0$, with equality if and only if $q=g$, the unique closest Gaussian to $p$ in forward KL is the moment-matched Gaussian $g$. The gap equals $h(g)-h(p)$, so $h(p)\le h(g)$ with equality if and only if $p=g$.

\section{Autoregressive factorization of a nondegenerate Gaussian}

Partition $\Sigma\succ 0$ so that the last coordinate is $x_d$. The Gaussian conditional of Section~3 gives
\[
x_d\given x_{<d}\sim\N(w_d^\top x_{<d},\sigma_d^2),
\]
with unique coefficients $(w_d,\sigma_d^2)$, and the $d=1$ factor equal to $\N(0,\Sigma_{11})$. Hence
\[
p(x)=\prod_{d=1}^D \N(x_d\given w_d^\top x_{<d},\sigma_d^2).
\]
The conditional variance is the Schur complement of the leading principal submatrix, so $\sigma_d^2=\det\Sigma_{1:d}/\det\Sigma_{1:d-1}$ and
\begin{resultbox}[Determinant factorization]
\[
\det\Sigma=\prod_{d=1}^D\sigma_d^2.
\]
\end{resultbox}
A strictly lower-triangular $W$ has $D(D-1)/2$ entries and there are $D$ positive conditional variances, matching the $D(D+1)/2$ free parameters of $\Sigma\succ 0$. Conversely, $x_d=w_d^\top x_{<d}+\varepsilon_d$ with independent $\varepsilon_d\sim\N(0,\sigma_d^2)$ yields
\[
x\sim\N\bigl(0,\;(I-W)^{-1}\Delta(I-W)^{-\top}\bigr),\qquad\Delta=\diag(\sigma_d^2)\succ 0.
\]
The sparsity pattern of $W$ depends on the coordinate order: a hub-and-leaves covariance is sparse in the hub-first order and dense in the leaf-first order. Heteroscedastic conditionals such as $\mathrm{Var}(x_2\given x_1)=e^{ux_1}$ cannot arise from a jointly nondegenerate Gaussian, because the conditional covariance of a Gaussian never depends on the value of the conditioner.

\section{Gaussian conditional negative log-likelihood}

An autoregressive continuous model often emits, at each coordinate, a univariate Gaussian with mean $\mu$ and log-variance $s$, so the variance is $e^s$. The negative log-density is
\begin{resultbox}[Univariate Gaussian NLL]
\begin{equation}
\label{eq:gauss-nll}
\ell(x;\mu,s)=\tfrac12\log(2\pi)+\tfrac s2+\frac{(x-\mu)^2}{2e^s}.
\end{equation}
\end{resultbox}
The gradients are $\partial\ell/\partial\mu=(\mu-x)e^{-s}$ and $\partial\ell/\partial s=\tfrac12-(x-\mu)^2/(2e^s)$. The Hessian in $(\mu,s)$ has determinant $-\tfrac12(x-\mu)^2 e^{-2s}\le 0$, so the loss is not jointly convex in $(\mu,s)$. In the natural parameters $\eta_1=\mu e^{-s}$ and $\eta_2=-\tfrac12 e^{-s}$ the same loss is strictly convex, because it is an exponential-family negative log-likelihood. Separate convexity in $\mu$ (for fixed $s$) and in $s$ (for fixed $\mu$) still holds.

A single constant Gaussian fit of $N$ samples has critical point $\hat\mu=\bar x$ and $e^{\hat s}=N^{-1}\sum(x_n-\bar x)^2$ whenever the sample variance is positive. With one training point the NLL is unbounded below ($s\to-\infty$ with $\mu=x$). With two or more distinct points it is bounded below by $\tfrac N2\log(2\pi e S/N)$. A flexible conditional model that can interpolate $(\mu_d,s_d)$ at the training inputs can again send $s\to-\infty$ and drive the training NLL to $-\infty$; a hard floor $s_d\ge s_{\min}$ restores a finite lower bound. Mixtures of Gaussians inherit the same escape: one component collapses onto a single observation while another keeps the remaining points at finite density.

\section{Homework 1 problems}

\begin{hwqbox}[Homework 1, problem 1.1]
Let $x\sim\N(\mu,\Sigma)$ and $y=Ax+b$. Using the moment generating function, derive the law of $y$, allowing singular $A\Sigma A^\top$.
\end{hwqbox}
\begin{hwsolbox}[Homework 1, problem 1.1]
For every $s$,
\begin{align*}
\E[e^{s^\top y}]
&=\E[e^{s^\top(Ax+b)}]
=e^{s^\top b}\,\E[e^{(A^\top s)^\top x}]\\
&=e^{s^\top b}\exp\bigl((A^\top s)^\top\mu+\tfrac12(A^\top s)^\top\Sigma A^\top s\bigr)\\
&=\exp\bigl(s^\top(A\mu+b)+\tfrac12 s^\top(A\Sigma A^\top)s\bigr).
\end{align*}
Hence $y\sim\N(A\mu+b,\,A\Sigma A^\top)$, possibly degenerate.
\end{hwsolbox}

\begin{hwqbox}[Homework 1, problem 1.1(a)]
Let $a,b\in\R$ and let $(x,y)$ be jointly Gaussian with cross-covariance $C=\mathrm{Cov}(x,y)$. Derive the law of $z=ax+by$. Discuss singularity when $x\perp y$ or $y=x$, and deduce $|\rho|\le 1$ in dimension one.
\end{hwqbox}
\begin{hwsolbox}[Homework 1, problem 1.1(a)]
\[
z=ax+by\sim\N\bigl(a\mu_x+b\mu_y,\; a^2\Sigma_x+b^2\Sigma_y+ab\,C+ab\,C^\top\bigr),
\]
because $\mathrm{Cov}(ax,by)=ab\,C$.
\begin{itemize}
\item If $x\perp y$, then $C=0$. If $\Sigma_x,\Sigma_y\succ 0$ and $(a,b)\ne(0,0)$, the covariance $a^2\Sigma_x+b^2\Sigma_y$ is positive definite, so not singular.
\item If $y=x$, then $z=(a+b)x$ and the covariance is $(a+b)^2\Sigma_x$. It is singular whenever $a+b=0$ and also whenever $\Sigma_x$ itself is singular.
\end{itemize}
In dimension $d=1$, with $\sigma_x,\sigma_y>0$ and correlation $\rho=C/(\sigma_x\sigma_y)$,
\[
\mathrm{Var}(z)=a^2\sigma_x^2+b^2\sigma_y^2+2ab\rho\,\sigma_x\sigma_y.
\]
The right-hand side is a nonnegative quadratic form in $(a,b)$ for every $(a,b)$, so its matrix has determinant $\sigma_x^2\sigma_y^2(1-\rho^2)\ge 0$. Thus $|\rho|\le 1$. Equality $|\rho|=1$ holds if and only if that matrix is singular, if and only if $y-\mu_y$ is an affine function of $x-\mu_x$ almost surely. This is Cauchy--Schwarz on the centered variables.
\end{hwsolbox}

\begin{hwqbox}[Homework 1, problem 1.1(b)]
Let $x\sim\N(0,I_d)$ and $A\in\R^{m\times d}$ with $m>d$ and full column rank. Show that $Ax$ has no density on $\R^m$. Using a thin QR factorization $A=QR$, relate the second moment $\E[\lVert Ax\rVert^2]$ to $\lVert A\rVert_F^2$.
\end{hwqbox}
\begin{hwsolbox}[Homework 1, problem 1.1(b)]
$\mathrm{Cov}(Ax)=AA^\top$ has rank $d<m$, so the law of $Ax$ is supported on the column space of $A$. It is not absolutely continuous on $\R^m$ and has no density there.

Thin QR: $A=QR$ with $Q\in\R^{m\times d}$, $Q^\top Q=I_d$, and $R\in\R^{d\times d}$ invertible. Then $Rx\sim\N(0,RR^\top)$ with $RR^\top\succ 0$, and $Ax=Qz$ with that nondegenerate $z\in\R^d$.

Two computations of the second moment:
\[
\E[\lVert Ax\rVert^2]=\E[x^\top A^\top Ax]=\tr(A^\top A)=\lVert A\rVert_F^2,
\]
and, since $Q$ has orthonormal columns, $\lVert Ax\rVert^2=\lVert z\rVert^2$ and $\E[\lVert z\rVert^2]=\tr(RR^\top)=\lVert R\rVert_F^2=\lVert A\rVert_F^2$.
\end{hwsolbox}

\begin{hwqbox}[Homework 1, problem 1.2]
Derive the product identity for two nondegenerate Gaussian densities, including the parallel-sum forms of $(c,C)$ and the normalizing constant $Z=\N(a\given b,A+B)$.
\end{hwqbox}
\begin{hwsolbox}[Homework 1, problem 1.2]
The sum of the two precision quadratics has precision $A^{-1}+B^{-1}$ and linear term $A^{-1}a+B^{-1}b$. Therefore
\[
C=(A^{-1}+B^{-1})^{-1},\qquad c=C(A^{-1}a+B^{-1}b).
\]
The parallel-sum identities $A^{-1}+B^{-1}=B^{-1}(A+B)A^{-1}$ give
\[
C=A(A+B)^{-1}B=B(A+B)^{-1}A,\qquad
c=B(A+B)^{-1}a+A(A+B)^{-1}b.
\]
They agree with the precision form because $CA^{-1}a=B(A+B)^{-1}a$ and $CB^{-1}b=A(A+B)^{-1}b$.

Integrating the product in $x$ produces $Z$, since $\N(x\given c,C)$ integrates to $1$. If $U\sim\N(a,A)$ and $V\sim\N(b,B)$ are independent, that integral is the density of $U-V$ at $0$, and $U-V\sim\N(a-b,A+B)$. Hence
\[
Z=\N(a\given b,\,A+B)=\N(b\given a,\,A+B).
\]
\end{hwsolbox}

\begin{hwqbox}[Homework 1, problem 1.2(a)]
With prior $\theta\sim\N(m_0,s_0^2)$ and i.i.d.\ observations $x_i\given\theta\sim\N(\theta,s^2)$, derive the Gaussian posterior, the marginal of the sample, the MLE, and the MAP. Discuss the limits $n\to\infty$ and $s_0\to\infty$.
\end{hwqbox}
\begin{hwsolbox}[Homework 1, problem 1.2(a)]
Likelihood in $\theta$: $\bar x\given\theta\sim\N(\theta,s^2/n)$, and $\N(\bar x\given\theta,s^2/n)=\N(\theta\given\bar x,s^2/n)$. Combine with the prior by the product identity:
\[
\frac1{s_n^2}=\frac1{s_0^2}+\frac n{s^2},\qquad
s_n^2=\frac{s_0^2 s^2}{s^2+n s_0^2},
\]
\[
m_n=s_n^2\Bigl(\frac{m_0}{s_0^2}+\frac{n\bar x}{s^2}\Bigr)
=\frac{s^2}{s^2+n s_0^2}\,m_0+\frac{n s_0^2}{s^2+n s_0^2}\,\bar x.
\]
So $\theta\given x_{1:n}\sim\N(m_n,s_n^2)$. The posterior mean is the stated mixture of the prior mean and the MLE.

Marginal of the sample: $\E[x_i]=m_0$, $\mathrm{Var}(x_i)=s^2+s_0^2$, and $\mathrm{Cov}(x_i,x_j)=s_0^2$ for $i\ne j$. Thus
\[
(x_1,\ldots,x_n)\sim\N\bigl(m_0\mathbf 1,\; s^2 I+s_0^2\mathbf 1\mathbf 1^\top\bigr).
\]
The likelihood $\prod_i\N(x_i\given\theta,s^2)$ is maximized at the MLE $\hat\theta_{\mathrm{MLE}}=\bar x$. The MAP is the posterior mode $m_n$.

As $n\to\infty$, the weight on $m_0$ tends to $0$, the weight on $\bar x$ tends to $1$, and $s_n^2\sim s^2/n\to 0$. The posterior collapses at the MLE. As $s_0\to\infty$, $s_n^2\to s^2/n$ and $m_n\to\bar x$: a flat prior makes MAP equal MLE.
\end{hwsolbox}

\begin{hwqbox}[Homework 1, problem 1.3]
Derive the conditional law of a Gaussian block $y$ given $x$ without inverting the full joint covariance. Show that the conditional covariance does not depend on the observation and that $\mathrm{Cov}(y\given x)\preceq\Sigma_{yy}$.
\end{hwqbox}
\begin{hwsolbox}[Homework 1, problem 1.3]
Set $K=\Sigma_{yx}\Sigma_{xx}^{-1}$, which exists because $\Sigma_{xx}\succ 0$. The residual
\[
r=y-\mu_y-K(x-\mu_x)
\]
is uncorrelated with $x$: $\mathrm{Cov}(r,x)=\Sigma_{yx}-K\Sigma_{xx}=0$. The pair $(r,x)$ is jointly Gaussian, so $r\perp x$. Therefore
\[
y\given x=x_0\sim\N\bigl(\mu_y+K(x_0-\mu_x),\;\mathrm{Cov}(r)\bigr),
\]
\[
\mathrm{Cov}(r)=\Sigma_{yy}-\Sigma_{yx}\Sigma_{xx}^{-1}\Sigma_{xy}.
\]
The conditional covariance does not depend on $x_0$. The subtracted term is positive semidefinite, so $\mathrm{Cov}(y\given x)\preceq\Sigma_{yy}$. Equality holds if and only if $\Sigma_{yx}=0$. For jointly Gaussian vectors that is exactly independence of $x$ and $y$.
\end{hwsolbox}

\begin{hwqbox}[Homework 1, problem 1.3(a)]
For the linear-Gaussian model $y=Ax+b+\varepsilon$ with $\varepsilon\sim\N(0,S)$ independent of $x\sim\N(\mu,\Sigma)$, derive both posterior forms for $x\given y$ and prove they agree via Woodbury's identity. Interpret the coefficient of the innovation.
\end{hwqbox}
\begin{hwsolbox}[Homework 1, problem 1.3(a)]
The joint is
\[
\begin{pmatrix}x\\ y\end{pmatrix}
\sim\N\Bigg(
\begin{pmatrix}\mu\\ A\mu+b\end{pmatrix},
\begin{pmatrix}\Sigma & \Sigma A^\top\\ A\Sigma & S+A\Sigma A^\top\end{pmatrix}
\Bigg).
\]
The marginal is $y\sim\N(A\mu+b,\,S+A\Sigma A^\top)$. Conditioning by problem~1.3 with $R:=S+A\Sigma A^\top$ gives
\[
\mathrm{Cov}(x\given y)=\Sigma-\Sigma A^\top R^{-1}A\Sigma,\qquad
\E[x\given y]=\mu+\Sigma A^\top R^{-1}\bigl(y-(A\mu+b)\bigr).
\]
The information form is
\[
\Sigma_{x\given y}=(\Sigma^{-1}+A^\top S^{-1}A)^{-1},\qquad
\E[x\given y]=\Sigma_{x\given y}\bigl(\Sigma^{-1}\mu+A^\top S^{-1}(y-b)\bigr).
\]
Woodbury's identity is exactly
\[
(\Sigma^{-1}+A^\top S^{-1}A)^{-1}=\Sigma-\Sigma A^\top R^{-1}A\Sigma,
\]
so the covariances agree. For the means, $\Sigma_{x\given y}\Sigma^{-1}=I-\Sigma_{x\given y}A^\top S^{-1}A$, hence the information-form mean equals
\[
\mu+\Sigma_{x\given y}A^\top S^{-1}\bigl(y-(A\mu+b)\bigr).
\]
The remaining identity $\Sigma_{x\given y}A^\top S^{-1}=\Sigma A^\top R^{-1}$ follows by multiplying on the left by $\Sigma^{-1}+A^\top S^{-1}A$:
\begin{align*}
(\Sigma^{-1}+A^\top S^{-1}A)\Sigma A^\top R^{-1}
&=A^\top(I+S^{-1}A\Sigma A^\top)R^{-1}
=A^\top S^{-1}(S+A\Sigma A^\top)R^{-1}
=A^\top S^{-1}.
\end{align*}
The matrix in front of the innovation is $\mathrm{Cov}(x,y)\mathrm{Cov}(y)^{-1}$, the linear regression coefficient of $x$ on $y$.
\end{hwsolbox}

\begin{hwqbox}[Homework 1, problem 1.4]
For a Gaussian mixture with weights $\pi_k$ and components $\N(\mu_k,\Sigma_k)$, derive the mean and the within/between decomposition of the covariance.
\end{hwqbox}
\begin{hwsolbox}[Homework 1, problem 1.4]
Let $K$ be a discrete component with $\mathbb P(K=k)=\pi_k$, and $x\given K=k\sim\N(\mu_k,\Sigma_k)$. Then
\[
\E[x]=\sum_k\pi_k\mu_k=:\mu,\qquad
\E[xx^\top]=\sum_k\pi_k(\Sigma_k+\mu_k\mu_k^\top).
\]
Therefore
\[
\mathrm{Cov}(x)
=\underbrace{\sum_k\pi_k\Sigma_k}_{\text{within}}
+\underbrace{\sum_k\pi_k(\mu_k-\mu)(\mu_k-\mu)^\top}_{\text{between}}.
\]
\end{hwsolbox}

\begin{hwqbox}[Homework 1, problem 1.4(a)]
For the equal mixture of $\N(-a,1)$ and $\N(a,1)$ on the line, compute mean, variance, and excess kurtosis. Show that $p$ is not Gaussian. Determine the number of modes as a function of $a>0$.
\end{hwqbox}
\begin{hwsolbox}[Homework 1, problem 1.4(a)]
Mean: $\tfrac12(-a)+\tfrac12 a=0$. Variance: within contribution $1$, between contribution $a^2$, so $\mathrm{Var}(x)=1+a^2$.

For $Z\sim\N(a,1)$, $\E[Z^4]=a^4+6a^2+3$. The same value holds at $-a$, so $\E_p[x^4]=a^4+6a^2+3$. Excess kurtosis:
\[
\frac{\E[x^4]}{\mathrm{Var}(x)^2}-3
=\frac{a^4+6a^2+3}{(1+a^2)^2}-3
=\frac{-2a^4}{(1+a^2)^2}<0
\quad(a>0).
\]
A nondegenerate Gaussian has excess kurtosis $0$, so $p$ is not Gaussian. Equivalently, its mgf is $\cosh(at)\,e^{t^2/2}$, which is not the mgf of a Gaussian.

Modes. The unnormalized density $h(x)=e^{-(x+a)^2/2}+e^{-(x-a)^2/2}$ has every critical point in $(-a,a)$, where
\[
h'(x)=0\iff \frac{a+x}{a-x}=e^{2ax}.
\]
The function $k(x)=\log(a+x)-\log(a-x)-2ax$ satisfies $k(0)=0$ and
\[
k'(x)=2a\frac{x^2-(a^2-1)}{a^2-x^2}.
\]
Also $h''(0)=2(a^2-1)e^{-a^2/2}$.
\begin{itemize}
\item If $a<1$, then $k'(x)>0$ on $(0,a)$, so $x=0$ is the only root and $h''(0)<0$: one mode.
\item If $a=1$, still only one root; near $0$, $h''(x)=-2x^2+O(x^4)$, so $x=0$ is a strict local maximum: one mode.
\item If $a>1$, then $k'$ changes from negative to positive at $\sqrt{a^2-1}$, and $k(a^-)=+\infty$. So $k$ has exactly one positive root. Together with $x=0$ and the negative root there are three critical points, with $h''(0)>0$: two modes.
\end{itemize}
Thus $p$ has two modes if and only if $a>1$.
\end{hwsolbox}

\begin{hwqbox}[Homework 1, problem 1.4(b)]
(i) Derive the multivariate Gaussian KL formula~\eqref{eq:mvnkl}. (ii) Prove the Pythagorean identity $\KL(p\Vert q)=\KL(p\Vert g)+\KL(g\Vert q)$ for $g=\N(\mu_p,\Sigma_p)$ and any Gaussian $q$, deduce the maximum-entropy property of Gaussians, and apply it to the mixture in (a).
\end{hwqbox}
\begin{hwsolbox}[Homework 1, problem 1.4(b)]
(i) For $p_i=\N(\mu_i,\Sigma_i)$, $\Sigma_i\succ 0$, formula~\eqref{eq:mvnkl} holds. Proof: $-\E_{p_0}\log p_1$ contributes $\tfrac12\log\det\Sigma_1+\tfrac12\E[(x-\mu_1)^\top\Sigma_1^{-1}(x-\mu_1)]$, and
\[
\E_{p_0}[(x-\mu_1)^\top\Sigma_1^{-1}(x-\mu_1)]
=\tr(\Sigma_1^{-1}\Sigma_0)+(\mu_0-\mu_1)^\top\Sigma_1^{-1}(\mu_0-\mu_1).
\]
Subtracting $h(p_0)=\tfrac d2\log(2\pi e)+\tfrac12\log\det\Sigma_0$, and cancelling $\tfrac d2\log(2\pi)$, leaves~\eqref{eq:mvnkl}.

(ii) Let $g=\N(\mu_p,\Sigma_p)$ and let $q$ be any Gaussian. The function $\log(g/q)$ is a polynomial of degree at most $2$. Any such polynomial has the same expectation under $p$ and under $g$, because $p$ and $g$ share their moments of order $0,1,2$. Therefore
\[
\KL(p\Vert q)-\KL(p\Vert g)
=\E_p\log\frac gq
=\E_g\log\frac gq
=\KL(g\Vert q),
\]
which is $\KL(p\Vert q)=\KL(p\Vert g)+\KL(g\Vert q)$. Since $\KL(g\Vert q)\ge 0$, with equality if and only if $q=g$, the unique minimizer of $\KL(p\Vert q)$ over Gaussians $q$ is $q^\star=g$. The minimum value is $\KL(p\Vert g)=h(g)-h(p)$, because $\E_p\log g=\E_g\log g=-h(g)$. Applying (i) to the pair $(g,q)$ writes the excess $\KL(g\Vert q)$ in closed form.

In particular $h(p)\le h(g)$, with equality if and only if $p=g$. Among densities with a fixed mean and a fixed positive definite covariance, the Gaussian has the largest differential entropy.

Applied to the mixture in (a): the moment-matched Gaussian is $\N(0,1+a^2)$, and the minimal KL to a Gaussian is $h(\N(0,1+a^2))-h(p)>0$.
\end{hwsolbox}

\begin{hwqbox}[Homework 1, problem 1.5]
Show that every nondegenerate Gaussian on $\R^D$ factors uniquely as $\prod_d\N(x_d\given w_d^\top x_{<d},\sigma_d^2)$, prove $\det\Sigma=\prod_d\sigma_d^2$, and recover the covariance as $(I-W)^{-1}\Delta(I-W)^{-\top}$.
\end{hwqbox}
\begin{hwsolbox}[Homework 1, problem 1.5]
Partition $\Sigma$ so that the last coordinate is $x_d$. By problem~1.3, with $x_{<d}$ in the role of the conditioning variable,
\[
w_d=\Sigma_{<d}^{-1}\Sigma_{<d,d},\qquad
\sigma_d^2=\Sigma_{dd}-\Sigma_{d,<d}\Sigma_{<d}^{-1}\Sigma_{<d,d},
\]
and $x_d\given x_{<d}\sim\N(w_d^\top x_{<d},\sigma_d^2)$, with the $d=1$ factor equal to $\N(0,\Sigma_{11})$. The conditional law of a Gaussian is unique, so $(w_d,\sigma_d^2)$ is unique, and so is the factorization.

The conditional variance is the Schur complement of the leading principal submatrix $\Sigma_{1:d-1}$. The determinant factorization of a Schur complement gives $\sigma_d^2=\det\Sigma_{1:d}/\det\Sigma_{1:d-1}$, with the convention $\det\Sigma_{1:0}=1$. Therefore $\det\Sigma=\prod_{d=1}^D\sigma_d^2$.

A strictly lower-triangular $W$ has $D(D-1)/2$ entries and there are $D$ positive conditional variances, hence $D(D+1)/2$ free parameters, the same number as a positive definite covariance. Every $\Sigma\succ 0$ arises exactly once.

Conversely, $x_d=w_d^\top x_{<d}+\varepsilon_d$ with independent $\varepsilon_d\sim\N(0,\sigma_d^2)$ is $x=(I-W)^{-1}\varepsilon$. The matrix $I-W$ is unit lower triangular, hence invertible, and $\varepsilon\sim\N(0,\diag(\sigma_d^2))$ is nondegenerate. Thus
\[
x\sim\N\big(0,\;(I-W)^{-1}\Delta(I-W)^{-\top}\big),\qquad\Delta=\diag(\sigma_d^2)\succ 0.
\]
\end{hwsolbox}

\begin{hwqbox}[Homework 1, problem 1.5(a)]
Identify the autoregressive factorization of the stationary AR(1) covariance $\Sigma_{ij}=\sigma^2\rho^{|i-j|}/(1-\rho^2)$.
\end{hwqbox}
\begin{hwsolbox}[Homework 1, problem 1.5(a)]
The matrix is the covariance of
\[
x_1\sim\N\Big(0,\frac{\sigma^2}{1-\rho^2}\Big),\qquad
x_d=\rho x_{d-1}+\varepsilon_d,\quad\varepsilon_d\sim\N(0,\sigma^2)\ \text{i.i.d.}
\]
For $d=1$, $\sigma_1^2=\sigma^2/(1-\rho^2)$. For $d\ge 2$, the candidate $w_d=\rho\,e_{d-1}$ satisfies the normal equations because
\[
\rho\cdot\frac{\sigma^2}{1-\rho^2}\rho^{|(j)-(d-1)|}
=\frac{\sigma^2}{1-\rho^2}\rho^{|d-j|}.
\]
The conditional variance collapses to $\sigma_d^2=\sigma^2$. So $w_d$ has a single nonzero entry, equal to $\rho$, in the coordinate of $x_{d-1}$.
\end{hwsolbox}

\begin{hwqbox}[Homework 1, problem 1.5(b)]
For a hub-and-leaves Gaussian with hub variance $1$, leaf residual variance $s^2$, and hub--leaf covariance $\beta$, compute the autoregressive factorization in hub-first and leaf-first orders. Show that the product of conditional variances equals $\det\Sigma$ in both orders, and that the hub-first sparsity cannot be transplanted to the reversed order when $D\ge 3$.
\end{hwqbox}
\begin{hwsolbox}[Homework 1, problem 1.5(b)]
In blocks,
\[
\Sigma=\begin{pmatrix}1 & \beta\mathbf 1^\top\\ \beta\mathbf 1 & s^2 I+\beta^2\mathbf 1\mathbf 1^\top\end{pmatrix}.
\]
\textbf{Order $(1,\ldots,D)$.} Here $\sigma_1^2=1$. For each leaf, $x_j=\beta x_1+\varepsilon_j$ with $\varepsilon_j$ independent of $(x_1,\ldots,x_{j-1})$, so $w_j=\beta\,e_1$ and $\sigma_j^2=s^2$. There are $D-1$ nonzero weights, and $\prod_d\sigma_d^2=(s^2)^{D-1}$.

\textbf{Order $(D,\ldots,1)$.} Position $m+1$ ($0\le m\le D-2$) is a leaf conditioned on $m$ other leaves. The constant weight $c=\beta^2/(s^2+m\beta^2)$ solves the normal equations, and
\[
\sigma^2
=\beta^2+s^2-\frac{m\beta^4}{s^2+m\beta^2}
=\frac{s^2\big(s^2+(m+1)\beta^2\big)}{s^2+m\beta^2}.
\]
The hub is last, with weight $\beta/(s^2+(D-1)\beta^2)$ on every leaf and conditional variance $\sigma_{\mathrm{hub}}^2=s^2/(s^2+(D-1)\beta^2)$. Nonzero weights: $\sum_{m=1}^{D-2}m+(D-1)=D(D-1)/2$.

The product of conditional variances equals $\det\Sigma$ in every order. Multiplying the leaf formulas telescopes to $(s^2)^{D-2}(s^2+(D-1)\beta^2)$, and the hub factor cancels the second piece, leaving $(s^2)^{D-1}$.

Sparsity cannot be transplanted: for a fixed ordering the pair $(W,\Delta)$ is unique. In the leaf-first ordering, the row in position $m+1$ has $m$ nonzero entries. For $D\ge 3$ and $m\ge 2$ that is already more than one, whereas every row of the hub-first factorization has at most one nonzero weight.
\end{hwsolbox}

\begin{hwqbox}[Homework 1, problem 1.5(c)]
Can the heteroscedastic conditional $x_1\sim\N(0,1)$, $x_2\given x_1\sim\N(0,e^{ux_1})$ with $u\ne 0$ arise from a jointly nondegenerate Gaussian? Justify.
\end{hwqbox}
\begin{hwsolbox}[Homework 1, problem 1.5(c)]
No. Under a jointly nondegenerate Gaussian, problem~1.3 says that $\mathrm{Cov}(x_2\given x_1)$ does not depend on the value of $x_1$. Here that conditional variance is $e^{ux_1}$, which is nonconstant because $u\ne 0$.

As a second witness: $\E[x_2^4]=3e^{2u^2}$ and $(\E[x_2^2])^2=e^{u^2}$, so the excess kurtosis is $3(e^{u^2}-1)>0$.
\end{hwsolbox}

\begin{hwqbox}[Homework 1, problem 5.1]
For the Gaussian NLL~\eqref{eq:gauss-nll}, compute $\partial\ell/\partial\mu$ and $\partial\ell/\partial s$. For an autoregressive model with $\mu_1=b_1$, $s_1=e_1$, and $\mu_d=v_d^\top h_d+b_d$, $s_d=u_d^\top h_d+e_d$ for $d\ge 2$, derive the gradients with respect to the biases, output weights, hidden state, pre-activations, and shared parameters $(c,W)$.
\end{hwqbox}
\begin{hwsolbox}[Homework 1, problem 5.1]
\[
\delta^\mu:=\frac{\partial\ell}{\partial\mu}=(\mu-x)e^{-s},\qquad
\delta^s:=\frac{\partial\ell}{\partial s}=\tfrac12-\frac{(x-\mu)^2}{2e^s}.
\]
Since $L=\sum_d\ell(x_d;\mu_d,s_d)$,
\begin{align*}
\frac{\partial L}{\partial b_d}&=\delta^\mu_d, &
\frac{\partial L}{\partial e_d}&=\delta^s_d,\\
\frac{\partial L}{\partial v_d}&=\delta^\mu_d\,h_d\ (d\ge 2), &
\frac{\partial L}{\partial u_d}&=\delta^s_d\,h_d\ (d\ge 2),\\
\frac{\partial L}{\partial h_d}&=\delta^\mu_d v_d+\delta^s_d u_d, &
\frac{\partial L}{\partial a_d}&=\phi'(a_d)\odot\frac{\partial L}{\partial h_d}.
\end{align*}
The shared parameters accumulate as
\[
\frac{\partial L}{\partial c}=\sum_{d\ge 2}g_d,\qquad
\frac{\partial L}{\partial W_{:,k}}=x_k\sum_{d>k}g_d,
\]
with $g_d=\partial L/\partial a_d$.
\end{hwsolbox}

\begin{hwqbox}[Homework 1, problem 5.1(a)]
For a constant Gaussian fit of $N$ i.i.d.\ samples, find the critical point in $(\mu,s)$, argue it is the global minimum when the sample variance is positive, and compute $\E[e^{\hat s}]$ under $x_n\sim\N(\mu_0,\sigma_0^2)$.
\end{hwqbox}
\begin{hwsolbox}[Homework 1, problem 5.1(a)]
The stationarity condition $\sum_n(\mu-x_n)e^{-s}=0$ forces $\mu=\bar x$. Then $\sum_n\delta^s=0$ forces
\[
e^{\hat s}=\frac1N\sum_n(x_n-\bar x)^2.
\]
If the sample variance is positive, this is the only critical point. It is the global minimum: for any $(\mu,s)$,
\[
\sum_n\ell(x_n;\mu,s)\ge \tfrac N2\log(2\pi)+\tfrac N2 s+\frac{S_N}{2e^s},
\]
where $S_N=\sum(x_n-\bar x)^2\le\sum(x_n-\mu)^2$, and the right-hand side, as a function of $t=e^s$, is minimized at $t=S_N/N$. As $|\mu|\to\infty$ or $s\to\pm\infty$ the loss tends to infinity whenever $S_N>0$. If every $x_n$ equals the same value, then $s\to-\infty$ with $\mu$ equal to that value sends the loss to $-\infty$.

At the critical point, $\sum(x_n-\bar x)^2/\sigma_0^2\sim\chi^2_{N-1}$, so $\E[e^{\hat s}]=\sigma_0^2(N-1)/N$.
\end{hwsolbox}

\begin{hwqbox}[Homework 1, problem 5.2]
Compute the Hessian of $\ell$ in $(\mu,s)$ and show that it is indefinite whenever $x\ne\mu$. Conclude that the NLL is not jointly convex (nor jointly concave) in $(\mu,s)$.
\end{hwqbox}
\begin{hwsolbox}[Homework 1, problem 5.2]
\[
\nabla^2\ell=\begin{pmatrix}
e^{-s} & (x-\mu)e^{-s}\\
(x-\mu)e^{-s} & \tfrac12(x-\mu)^2 e^{-s}
\end{pmatrix}.
\]
Its determinant equals $-\tfrac12(x-\mu)^2 e^{-2s}\le 0$, and it is strictly negative whenever $x\ne\mu$. The Hessian is then indefinite. The loss is not jointly convex in $(\mu,s)$, and not jointly concave either.
\end{hwsolbox}

\begin{hwqbox}[Homework 1, problem 5.2(a)]
Rewrite the Gaussian density in natural parameters $\eta_1=\mu e^{-s}$, $\eta_2=-\tfrac12 e^{-s}$, and show that the NLL is strictly convex in $\eta$.
\end{hwqbox}
\begin{hwsolbox}[Homework 1, problem 5.2(a)]
Completing the exponential-family representation,
\[
\N(x\given\mu,e^s)=\frac{1}{\sqrt{2\pi}}\exp(\eta_1 x+\eta_2 x^2-A(\eta)),
\]
\[
A(\eta)=-\tfrac12\log(-2\eta_2)-\frac{\eta_1^2}{4\eta_2}.
\]
Hence $\ell=\tfrac12\log(2\pi)-\eta_1 x-\eta_2 x^2+A(\eta)$. The Hessian of $A$ is positive definite on $\{\eta_2<0\}$:
\[
\partial_{\eta_1\eta_1}A=\frac{1}{-2\eta_2}>0,\qquad
\det\nabla^2 A=-\frac{1}{4\eta_2^3}>0.
\]
So $A$, and therefore $\ell$, is strictly convex in $\eta$.
\end{hwsolbox}

\begin{hwqbox}[Homework 1, problem 5.2(b)]
Show separate convexity of $\ell$ in $\mu$ and in $s$. For network parameters, show convexity in the mean weights $(v_d,b_d)$ and in the variance weights $(u_d,e_d)$ separately, and explain why joint convexity in $(v_d,u_d)$ fails.
\end{hwqbox}
\begin{hwsolbox}[Homework 1, problem 5.2(b)]
For fixed $s$, $\partial^2\ell/\partial\mu^2=e^{-s}>0$: strictly convex in $\mu$. For fixed $\mu$, $\partial^2\ell/\partial s^2=\tfrac12(x-\mu)^2 e^{-s}\ge 0$: convex in $s$, strictly if $x\ne\mu$.

If $(u_d,e_d,h_d)$ are held fixed, $\mu_d$ is affine in $(v_d,b_d)$. A convex function of an affine map is convex, so the NLL is convex in $(v_d,b_d)$, jointly over $d$. Likewise, with everything else fixed, $s_d$ is affine in $(u_d,e_d)$ and the NLL is convex in those variance parameters.

It is not jointly convex in $(v_d,u_d)$. If $h_d\ne 0$, the map $(v_d,u_d)\mapsto(\mu_d,s_d)$ is linear and surjective onto $\R^2$. Joint convexity would imply joint convexity in $(\mu,s)$, which problem~5.2 forbids.
\end{hwsolbox}

\begin{hwqbox}[Homework 1, problem 5.3]
In one dimension, is the training NLL bounded below for a single training point? For $N\ge 2$ distinct points? Give the greatest lower bound in the latter case.
\end{hwqbox}
\begin{hwsolbox}[Homework 1, problem 5.3]
\textbf{One point.} Set $\mu=x$ and let $s\to-\infty$. Then $\ell=s/2+\tfrac12\log(2\pi)\to-\infty$. The NLL is not bounded below.

\textbf{$N\ge 2$ distinct points.} Let $S=\sum_n(x_n-\bar x)^2>0$. The inequality in problem~5.1(a) gives, for every $(\mu,s)$,
\[
\sum_{n=1}^N\ell(x_n;\mu,s)\ge\frac N2\log\Big(2\pi e\frac S N\Big),
\]
with equality at $\mu=\bar x$ and $e^s=S/N$. The average NLL is bounded below by $\tfrac12\log(2\pi e S/N)$.
\end{hwsolbox}

\begin{hwqbox}[Homework 1, problem 5.3(a)]
For a two-dimensional autoregressive Gaussian conditional model that can interpolate at the training inputs, is the training NLL bounded below? Contrast with a discrete model.
\end{hwqbox}
\begin{hwsolbox}[Homework 1, problem 5.3(a)]
No, the training NLL is not bounded below. The first coordinates are distinct, and the hidden layer is wide enough to interpolate arbitrary values of $(\mu_2,s_2)$ at those finitely many inputs. Set $\mu_2(x_{n1})=x_{n2}$ and send $s_2(x_{n1})\to-\infty$ at each training input, keeping $s_1$ fixed. Every conditional term $(x_{n2}-\mu_2)^2 e^{-s_2}$ vanishes while $s_2/2\to-\infty$, so the sum of NLLs tends to $-\infty$.

A discrete model cannot do the analogous thing. For any pmf $q$ and any data distribution $p_{\mathrm{data}}$,
\[
-\E_{p_{\mathrm{data}}}\log q=H(p_{\mathrm{data}})+\KL(p_{\mathrm{data}}\Vert q)\ge H(p_{\mathrm{data}}).
\]
The average NLL cannot undercut the entropy of the data distribution. Continuous training data are a sum of atoms; a density value at a point can be arbitrarily large, so there is no entropy floor under the empirical measure.
\end{hwsolbox}

\begin{hwqbox}[Homework 1, problem 5.3(b)]
Impose a variance floor $s_d\ge s_{\min}$. Show that the training NLL becomes bounded below, and identify what changed relative to problem~5.3(a).
\end{hwqbox}
\begin{hwsolbox}[Homework 1, problem 5.3(b)]
With $s_d\ge s_{\min}$ for every conditional,
\[
\ell(x_d;\mu_d,s_d)\ge\tfrac12\log(2\pi)+\tfrac{s_d}2\ge\tfrac12\big(\log(2\pi)+s_{\min}\big),
\]
because the squared residual is nonnegative and $s/2$ is increasing. Summing over the $D$ coordinates,
\[
-\log p_\theta(x)\ge\frac D2\big(\log(2\pi)+s_{\min}\big)
\]
for every $x$ and every parameter obeying the floor. The training NLL of problem~5.3(a) is therefore bounded below, by $N$ times this constant. The floor is what creates the bound: without it, $s_2\to-\infty$ was available at the training inputs.
\end{hwsolbox}

\begin{hwqbox}[Homework 1, problem 5.4(a)]
For a mixture of $C$ univariate Gaussians with mixture logits $\ell_c$, means $\mu_c$, and log-variances $s_c$, derive $\partial L/\partial\ell_c$, $\partial L/\partial\mu_c$, $\partial L/\partial s_c$, and $\partial L/\partial h$ when each natural parameter is linear in a shared hidden state $h$.
\end{hwqbox}
\begin{hwsolbox}[Homework 1, problem 5.4(a)]
Let $\varphi_c=\N(x\given\mu_c,e^{s_c})$ and $Z=\sum_j\pi_j\varphi_j$, so $L=-\log Z$ and $\gamma_c=\pi_c\varphi_c/Z$. Softmax Jacobian: $\partial\pi_j/\partial\ell_c=\pi_j(\mathbf 1_{j=c}-\pi_c)$. Hence
\[
\frac{\partial Z}{\partial\ell_c}=\pi_c\varphi_c-\pi_c Z,\qquad
\frac{\partial L}{\partial\ell_c}=\pi_c-\gamma_c.
\]
Summing over components, $\sum_c\partial L/\partial\ell_c=0$. Differentiating through one component,
\[
\frac{\partial L}{\partial\mu_c}=\gamma_c\,\delta^\mu_c,\qquad
\frac{\partial L}{\partial s_c}=\gamma_c\,\delta^s_c.
\]
Compared with problem~5.1, each Gaussian gradient is multiplied by the responsibility of that component. Since $\ell_c=a_c^\top h+\alpha_c$, $\mu_c=v_c^\top h+b_c$, and $s_c=u_c^\top h+e_c$,
\[
\frac{\partial L}{\partial h}
=\sum_{c=1}^C\Big((\pi_c-\gamma_c)a_c+\gamma_c\delta^\mu_c\,v_c+\gamma_c\delta^s_c\,u_c\Big).
\]
\end{hwsolbox}

\begin{hwqbox}[Homework 1, problem 5.4(b)]
For $C\ge 2$ and $N\ge 2$ distinct points, is the mixture likelihood bounded above? Contrast with $C=1$.
\end{hwqbox}
\begin{hwsolbox}[Homework 1, problem 5.4(b)]
No upper bound, once $C\ge 2$ and the $N\ge 2$ points are distinct. Put $\mu_1=x_1$, send $s_1\to-\infty$, and keep $\pi_1=\pi_2=1/2$, $\mu_2$ arbitrary, and $s_2$ fixed. The factor at $x_1$ tends to infinity because $\varphi_1(x_1)\to\infty$. At every $x_n\ne x_1$, $\varphi_1(x_n)\to 0$, but $\pi_2\varphi_2(x_n)$ stays equal to a positive constant. The product of the $N$ mixture densities tends to infinity. For $C=1$ this escape is impossible, by problem~5.3.
\end{hwsolbox}

\begin{hwqbox}[Homework 1, problem 5.4(c)]
Show that the mixture NLL is not convex, and identify the label-switching symmetry of order $C!$.
\end{hwqbox}
\begin{hwsolbox}[Homework 1, problem 5.4(c)]
Take $C=2$, one observation $x=0$, $\pi=(1/2,1/2)$, $s_1=s_2=0$, and free means. Up to an additive constant,
\[
L(\mu_1,\mu_2)=-\log\big(e^{-\mu_1^2/2}+e^{-\mu_2^2/2}\big).
\]
At the endpoints $(0,3)$ and $(3,0)$, $L=-\log(1+e^{-4.5})\approx -0.0111$. At the midpoint $(1.5,1.5)$, $L=-\log(2e^{-1.125})=-\log 2+1.125\approx 0.4319$, which is larger than the average of the endpoint values. The function lies above the chord, so it is not convex.

Swapping the two components' means leaves $(\pi,s)$ unchanged in this example and does not change the mixture. A generic pair with $\mu_1\ne\mu_2$ therefore has $2=2!$ parameter vectors, related by relabeling, that define the same distribution. For a mixture with $C$ labeled components the same relabeling symmetry has order $C!$.
\end{hwsolbox}
